% AI vs Numerical Analysis -- two forms of computation
% Claes Johnson. DRAFT. Companion to euler_bluff_body.tex
\documentclass[11pt]{article}
\usepackage[margin=1.15in]{geometry}
\usepackage{amsmath,amssymb}
\usepackage{url}
\usepackage[hidelinks]{hyperref}
\setlength{\parskip}{3pt}

\title{\vspace{-2em}AI versus Numerical Analysis:\\Two Forms of Computation}
\author{Claes Johnson\\[0.2em]
\small KTH Royal Institute of Technology, Stockholm\\
\small \texttt{claesjohnson@gmail.com}}
\date{\today}

\begin{document}
\maketitle

\begin{abstract}
\noindent
A large language model computes with symbols; numerical analysis computes with numbers in digital
form. Both are machine computation, both run on the same hardware, and the difference between them
is not rigour but what the computed objects denote and how the result is certified. A symbolic
computation is certified by derivability: each step follows, and a proof assistant can check that it
does. A numerical computation is certified by an error bound that the computation itself produces.
The first gives certainty about a statement; the second gives a quantity with a stated accuracy,
which can be compared with a measurement.

We take as example the two recent accounts of the incompressible Euler equations. An AI system has
produced a proof, formalised in Lean, that smooth compactly supported data can develop infinite
velocity gradients in finite time. Adaptive finite element computation produces turbulent solutions
past a body whose mean-value outputs --- drag, lift, separation --- are determined to a computed
tolerance and agree with measurement. The first proves \emph{existence without exhibiting an
individual case}; the second \emph{computes an individual case} and exhibits nothing general. Each
has exactly what the other lacks, neither is the other's refutation, and the pair is instructive
about both ways of computing.

We draw two conclusions. Symbolic computation reaches objects that possess exploitable structure ---
a symmetry, an ansatz, a closed form --- and reaches them exactly; it does not reach objects defined
by the absence of such structure, of which turbulence is the standing example, and no increase in
scale will alter this. And the problem AI now faces, that of trusting a computation no person has
followed, is one numerical analysis met fifty years ago and solved, by computing the error alongside
the answer. We suggest that a posteriori error control is what symbolic computation presently lacks
and what it should take from the older discipline.
\end{abstract}

\section{Two forms of computation}\label{sec:two}

A large language model and a finite element solver are both programs. Both execute on the same
processors, both reduce in the end to floating-point arithmetic, and both produce in finite time an
output no person could have produced by hand. The distinction between them is not that one is
mathematics and the other machinery.

The distinction is in what the computation manipulates and in what it means for the output to be
right.

\paragraph{Computation with symbols.} A language model operates on tokens --- words, formulas, proof
steps --- and its output is a text. When that text is a proof, correctness means \emph{derivability}:
every step follows from previous steps by a rule, and the question of whether it does is decidable
by a proof assistant. The verdict is binary. A proof is valid or it is not, and no quantity attaches
to it.

\paragraph{Computation with digital numbers.} A numerical method operates on numbers held in finite
precision over a finite mesh, and its output is a number. Correctness cannot mean derivability, since
the computed object is not the exact one and never will be. It means instead a \emph{bound}:
\[
|\,\text{output computed} - \text{output exact}\,| \;\le\; E ,
\]
where $E$ is itself computed, from the residual of the computed solution weighted by the solution of
an adjoint problem. The verdict is quantitative. The answer is right to within a stated amount, and
the amount is part of the answer.

\paragraph{Where the difference lies.} It is tempting to rank these, and the ranking usually proposed
is that derivability is the stronger notion because it is exact. That is true of the notion and
misleading about the practice. A derivation is exact about the statement it derives, and a statement
can be exactly established and say nothing about the world. A bound is inexact about a quantity, and
a quantity can be compared against a measurement made by somebody else. The two kinds of certainty
point in different directions: one inward, to consistency; one outward, to agreement.

\section{The example}\label{sec:example}

The incompressible Euler equations have lately received both treatments, which makes them an unusually
clean case.

\paragraph{The symbolic account.} An AI system has produced a proof that there exists smooth,
compactly supported, divergence-free initial data on $\mathbb{R}^3$ whose smooth Euler solution has
finite maximal lifespan: the velocity gradient becomes unbounded and the Beale--Kato--Majda vorticity
integral diverges~\cite{openaieuler}. The proof is fifty-seven pages and is formalised in Lean, so
that its correctness is machine-checked. Comparable results were being prepared independently and by
hand~\cite{buckmaster}.

What the result delivers is an existence statement. It does not say which data blow up; there is no
measure on the set and no claim of genericity. It does not say what the flow looks like. It yields
no drag, no lift, no length, no time that could be set beside an experiment. The datum itself is
produced as the limit of an iteration and cannot be written down. The result is nonetheless a real
advance: a question open for decades is settled, and the settling is exact.

\paragraph{The numerical account.} Adaptive finite element computation of the same equations past a
body produces turbulent solutions~\cite{book,dalembert,flight}. The velocity stays bounded, the
gradients do not, the solution dissipates energy at a rate independent of the viscosity, and the
limit of the computed sequence is a weak solution of Euler by the residual control the method
carries. Its mean-value outputs are determined: drag to within a few per cent of measurement, the
separation line where experiment puts it, the lift curve and the stall.

What this delivers is not a solution. For weak solutions of Euler uniqueness is false, so there is no
unique solution to deliver. What is delivered is the quantities that are determined, each with an
error bound the computation produces, each open to check by somebody with no interest in the outcome.

\paragraph{The shape of the comparison.} Put at its plainest, the difference is between a proof of
\emph{existence without an individual case} and the \emph{direct computation of an individual case}.

The symbolic account establishes that a datum with the stated property exists. It does not hand one
over. The datum is the limit of an iteration over stages, each adding an oscillation to the one
before, and no finite description of it is produced or intended --- nothing can be plotted, measured
or shown to anybody. What exists is the existence.

The numerical account computes one case: this body, this free stream, this flow, with this drag. It
establishes nothing about all data or about any other case. What it produces can be looked at. The
wake has a shape, the separation line a position, the drag a value with a bound on its error, and
each can be set beside a measurement.

So the two are complementary in the strict sense: each has exactly what the other lacks. Generality
without instance, instance without generality. The first is mathematically substantial and
physically empty, the second physically substantial and mathematically the weaker. The companion
paper~\cite{euler} argues the fluid mechanics of this and concludes that the singular solutions,
whatever their mathematical interest, are not flows. Here we are concerned with what the pair says
about the two ways of computing.

\section{What each form can reach}\label{sec:reach}

The division of labour is not a division between rigour and approximation. It follows from what each
form of computation requires of its object.

\paragraph{Symbolic computation requires exploitable structure.} A proof proceeds by rewriting, and
rewriting needs something to grip. Every construction of a singular solution buys its tractability
with a structure: self-similarity, which collapses a partial differential equation to a profile and
leaves a fixed point; or axisymmetry, which removes a variable; or an iterative scheme whose stages
are similar to one another. This is not a criticism. It is what makes the object reachable at all,
and the skill lies in finding the structure that a given problem will support.

\paragraph{Numerical computation requires none.} A mesh does not care whether the solution has a
symmetry. The computation proceeds identically whether the object is a self-similar profile or a
turbulent wake, and the cost is set by the resolution rather than by the structure.

\paragraph{Hence the complementarity, and its sharp edge.} Objects \emph{with} structure are reachable
both ways, and symbolically with greater precision. Objects \emph{without} structure are reachable
only numerically. Turbulence is the standing example of the second kind --- it is \emph{defined} by
carrying structure at every scale with no organising principle to exploit, so there is no ansatz to
be found, and the absence is the phenomenon rather than a gap in anybody's technique.

This bears directly on what AI changes. A language model is a search over symbol sequences, and what
it does superlatively is find structures that a human might have found given longer. It therefore
amplifies the half of mathematics that was already good at special objects, and amplifies it a great
deal. It does not convert the other half. No scale of model makes a turbulent flow analytically
constructible, because the obstruction is in the object and not in the searcher.

\section{Certainty, and what it is about}\label{sec:certainty}

Both accounts have now outrun human reading, and they have done so in the same week, which makes the
comparison of their epistemic positions unusually direct.

The proof is fifty-seven pages, machine-checked, and has not been verified in full by a person. The
computation produces an error bound no person derived by hand and which refers to an adjoint solution
no person inspected. In both cases a result is believed on the strength of a procedure rather than
of an understanding, and the discomfort this produces is being discussed as though it were new.

It is not new. It is the condition numerical analysis has worked in since the 1950s, and the
discipline's response was not to argue about whether machine results count. It was to make them
accountable: to compute, alongside the answer, a bound on how wrong the answer may be, and to refine
where the bound says the error lives. That is the content of a posteriori error estimation and
adaptivity, and it is why a computed drag coefficient is a scientific statement and not an opinion
produced by a program.

The two forms of certainty differ in what they are about, and the difference is worth stating
without rancour.

\emph{A formal proof is internally certain.} Lean guarantees that the steps follow. It guarantees
nothing about whether the statement proved is about anything, and in the present case the statement
is about a datum that cannot be prepared and a flow that no instrument will see. The certainty is
real and it is certainty of consistency.

\emph{A numerical computation is externally checkable.} Its output is a number with a tolerance, and
the number can be set beside a measurement that somebody else made, at a facility with no stake in
the result. The High Lift Prediction Workshop and the high-order workshops exist for exactly this:
predictions are submitted before the measurements are released. That is a form of accountability a
proof does not have and does not need, and a computation does need.

\section{What symbolic computation might take from the older discipline}\label{sec:take}

The central practical problem now facing AI in mathematics and elsewhere is this: a system produces
an output that is too long to check and too important to trust blindly. That is, exactly, the
problem numerical analysis faced when machines first outran hand calculation, and it has a
developed answer.

The answer has three parts, and all three transfer in principle.

\paragraph{Compute the error, not merely the answer.} A finite element computation returns a solution
and a bound on the error in each output of interest. The bound comes from the residual --- what the
computed solution fails to satisfy --- weighted by the solution of an adjoint problem which measures
how much that failure matters for the output in question. The structure is general: a residual, a
sensitivity, and their product. A language model has a residual in the same sense --- the respects
in which its output fails to satisfy the conditions it should --- and nothing at present weights it.

\paragraph{Let the error direct the effort.} Adaptivity refines where the bound says the error is,
not uniformly. The analogue for a symbolic system is to concentrate verification where the proof is
weakest rather than uniformly over its length, which requires knowing where that is, which requires
the weighting above.

\paragraph{Report what is \emph{not} determined.} An honest a posteriori estimate sometimes returns a
large bound, and that is information: the computation does not determine this output, and no
refinement of this kind will make it do so. A system that cannot return that verdict about its own
output is not yet accountable, whatever its success rate.

We put this forward as the constructive half of the comparison. The debate about whether machine
proof is mathematics will be settled the way the corresponding debate about machine computation was
settled, which is not by argument but by the development of error control. Numerical analysis has
fifty years of that, and it was built for precisely the situation symbolic computation is now in.

\section{Summary}\label{sec:summary}

Symbolic computation and numerical computation are both machine computation and differ in what they
manipulate and in how they are certified: derivability against an error bound, consistency against
agreement. Symbolic computation reaches objects possessing exploitable structure and reaches them
exactly; numerical computation reaches objects without it and reaches them to a stated tolerance.
Turbulence is of the second kind by definition, and no amount of symbolic power alters that. The
Euler equations have lately been treated both ways, and the results stand in exactly this relation:
a proof that blow-up data exists which exhibits no individual case, and the direct computation of an
individual case whose observables agree with measurement. Generality without instance, and instance
without generality.

The question now being asked of AI --- how to trust a result no one has followed --- is the question
numerical analysis answered by computing the error alongside the answer and refining where it lives.
It is the better part of what the older discipline has to offer the newer, and it is not yet taken up.

\begin{thebibliography}{99}

\bibitem{euler} J.~Hoffman and C.~Johnson, \emph{Incompressible Euler flow: turbulence vs blow-up},
companion paper, 2026.

\bibitem{openaieuler} OpenAI, \emph{Finite time blowup for the Euler equation}, preprint, 2026,
57~pp.
\newblock \url{https://cdn.openai.com/pdf/315b36cd-ec98-4023-8342-93345194ece1/euler.pdf}

\bibitem{buckmaster} T.~Buckmaster and L.~Alp\"oge, \emph{Statement}, 2026.
\newblock \url{https://cims.nyu.edu/~tristanb/statement.pdf}

\bibitem{book} J.~Hoffman and C.~Johnson, \emph{Computational Turbulent Incompressible Flow},
Applied Mathematics: Body and Soul, vol.~4, Springer, 2007.

\bibitem{dalembert} J.~Hoffman and C.~Johnson, \emph{Resolution of d'Alembert's paradox},
J.\ Math.\ Fluid Mech.\ \textbf{12} (2010) 321--334.

\bibitem{flight} J.~Hoffman, J.~Jansson and C.~Johnson, \emph{New theory of flight},
J.\ Math.\ Fluid Mech.\ \textbf{18} (2016) 219--241.

\bibitem{ek} K.~Eriksson, D.~Estep, P.~Hansbo and C.~Johnson, \emph{Computational Differential
Equations}, Cambridge University Press, 1996.

\bibitem{logg} A.~Logg, \emph{Automation of computational mathematical modeling}, Ph.D.\ thesis,
Chalmers University of Technology, 2004.

\bibitem{fenics} A.~Logg, K.-A.~Mardal and G.~N.~Wells (eds.), \emph{Automated Solution of
Differential Equations by the Finite Element Method}, Lecture Notes in Computational Science and
Engineering~84, Springer, 2012.

\end{thebibliography}

\end{document}
